\numoptstrue
\yearpaper{2016}{Version A-1 \textperiodcentered{} 60 Questions \textperiodcentered{} 60 Marks \textperiodcentered{} 80 Minutes}{}
\begin{multicols}{2}

\QQ{A body falls freely for $10\unit{sec}$. Its average velocity during this journey (take $g = 10\unit{m\,s^{-2}}$)\ntag{A}{17}{keph102}{5}{}{Class XI Part 1 \textperiodcentered{} Ch 2}\nsrc{None found in NEET/AIPMT/AIEEE/JEE papers, 2005--2026 (NCERT concept application)}}
{\opt{$100\unit{m\,s^{-1}}$}{$10\unit{m\,s^{-1}}$}{$50\unit{m\,s^{-1}}$}{$5\unit{m\,s^{-1}}$}}

\QQ{Three projectiles A, B and C are projected at an angle of $30^\circ$, $45^\circ$, $60^\circ$ respectively. If $R_A$, $R_B$ and $R_C$ are ranges of A, B and C respectively, then (velocity of projection is same for A, B \& C)\ntag{D}{40}{keph103}{14}{}{Class XI Part 1 \textperiodcentered{} Ch 3}}
{\opt{$R_A = R_B = R_C$}{$R_A = R_C > R_B$}{$R_A<R_B<R_C$}{$R_A = R_C < R_B$}}

\QQ{The component of a vector $\vec r$ along $x$-axis will have a maximum value if\ntag{D}{31}{keph103}{5}{}{Class XI Part 1 \textperiodcentered{} Ch 3}}
{\opt{$\vec r$ is along $+$ve $x$-axis}{$\vec r$ is along $+$ve $y$-axis}{$\vec r$ is along $-$ve $y$-axis}{$\vec r$ makes an angle of $45^\circ$ with the $x$-axis}}

\QQ{Maximum acceleration of the train in which a $50\unit{kg}$ box lying on its floor will remain stationary (Given: Co-efficient of static friction between the box and the train's floor is 0.3 and $g = 10\unit{m\,s^{-2}}$)\ntag{D}{61}{keph104}{13}{ \textperiodcentered{} Example 4.7}{Class XI Part 1 \textperiodcentered{} Ch 4}}
{\opt{$5.0\unit{m\,s^{-2}}$}{$3.0\unit{m\,s^{-2}}$}{$1.5\unit{m\,s^{-2}}$}{$15\unit{m\,s^{-2}}$}}

\QQ{A $12\unit{kg}$ bomb at rest explodes into two pieces of $4\unit{kg}$ and $8\unit{kg}$. If the momentum of $4\unit{kg}$ piece is $20\unit{N\,s}$, the kinetic energy of the $8\unit{kg}$ piece is\ntag{A}{57}{keph104}{9}{}{Class XI Part 1 \textperiodcentered{} Ch 4}\nsrc{None found in NEET/AIPMT/AIEEE/JEE papers, 2005--2026 (NCERT concept application)}}
{\opt{$25\unit{J}$}{$20\unit{J}$}{$50\unit{J}$}{$40\unit{J}$}}

\QQ{Which of the points is likely position of the centre of mass of the system shown in the figure?
\fig{\begin{tikzpicture}[scale=0.85]
\fill[pattern=north east lines,pattern color=gray] (0,0) -- (2.4,0) -- (2.4,1.6) -- cycle;
\draw (0,0) rectangle (2.4,1.6); \draw (0,0) -- (2.4,1.6);
\node[lbl,rotate=34] at (0.35,1.3) {Air}; \node[lbl,rotate=34] at (1.9,0.55) {Sand};
\foreach \p/\n in {(0.65,1.25)/A,(0.9,1.05)/B,(1.2,0.82)/C,(1.42,0.6)/D} {\fill \p circle (1.3pt); \node[font=\tiny,above right] at \p {\n};}
\end{tikzpicture}}\ntag{A}{98}{keph106}{7}{}{Class XI Part 1 \textperiodcentered{} Ch 6}\nsrc{None found in NEET/AIPMT/AIEEE/JEE papers, 2005--2026 (NCERT concept application)}}
{\opt{A}{D}{B}{C}}

\QQ{Three bodies a ring (R), a solid cylinder (C) and a solid sphere (S) having same mass and same radius roll down the inclined plane without slipping. They start from rest, if $v_R$, $v_C$ and $v_S$ are velocities of respective bodies on reaching the bottom of the plane, then\ntag{X}{}{}{}{}{rolling motion}}
{\opt{$v_R = v_C = v_S$}{$v_R>v_C>v_S$}{$v_R<v_C<v_S$}{$v_R = v_C > v_S$}}

\QQ{Variation of acceleration due to gravity ($g$) with distance $x$ from the centre of the earth is best represented by ($R\to$ Radius of the earth)\ntag{D}{134}{keph107}{8}{}{Class XI Part 1 \textperiodcentered{} Ch 7}}
{\optfig{\gaxes{$x$}{$g$}{\draw (0,1.5) -- (0.8,0.6) .. controls (1.1,1.1) and (1.4,1.25) .. (1.9,1.3); \draw[densely dotted] (0.8,0.6) -- (0.8,0) node[below,font=\tiny]{R};}}
{\gaxes{$x$}{$g$}{\draw (0,0) -- (1.6,1.7); \draw[densely dotted] (0.8,0.85) -- (0.8,0) node[below,font=\tiny]{R};}}
{\gaxes{$x$}{$g$}{\draw (0,1.6) -- (0.8,0.7) .. controls (1.1,0.35) and (1.4,0.25) .. (1.9,0.2); \draw[densely dotted] (0.8,0.7) -- (0.8,0) node[below,font=\tiny]{R};}}
{\gaxes{$x$}{$g$}{\draw (0,0) -- (0.8,1.4) .. controls (1.0,0.6) and (1.3,0.3) .. (1.9,0.2); \draw[densely dotted] (0.8,1.4) -- (0.8,0) node[below,font=\tiny]{R};}}}

\QQ{A spring is stretched by applying a load to its free end. The strain produced in the spring is\ntag{A}{170}{keph201}{4}{ \textperiodcentered{} Exemplar 9.4}{Class XI Part 2 \textperiodcentered{} Ch 8}\nsrc{None found in NEET/AIPMT/AIEEE/JEE papers, 2005--2026 (NCERT concept application)}}
{\opt{Volumetric}{Shear}{Longitudinal \& Shear}{Longitudinal}}

\QQ{An ideal fluid flows through a pipe of circular cross section with diameters $5\unit{cm}$ and $10\unit{cm}$ as shown. The ratio of velocities of fluid at A and B is
\fig{\begin{tikzpicture}[scale=0.8]
\draw (0,0.35) -- (2.6,0.7) (0,-0.35) -- (2.6,-0.7);
\draw (0,0) ellipse (0.12 and 0.35); \draw (2.6,0) ellipse (0.2 and 0.7);
\draw[->] (-0.5,0) -- (0.2,0); \draw[->] (2.6,0) -- (3.2,0);
\node[lbl,above] at (-0.3,0.3) {A}; \node[lbl,above right] at (2.6,0.6) {B};
\end{tikzpicture}}\ntag{D}{187}{keph202}{8}{}{Class XI Part 2 \textperiodcentered{} Ch 9}}
{\opt{4:1}{1:4}{2:1}{1:2}}

\QQ{A pan filled with hot food cools from $94^\circ$C to $86^\circ$C in 2 minutes. When the room temperature is $20^\circ$C. How long will it cool from $74^\circ$C to $66^\circ$C?\ntag{D}{221}{keph203}{20}{ \textperiodcentered{} Example 10.8}{Class XI Part 2 \textperiodcentered{} Ch 10}}
{\opt{2 minutes}{2.8 minutes}{2.5 minutes}{1.8 minutes}}

\QQ{Four rods with different radii $r$ and length $l$ are used to connect two heat reservoirs at different temperature. Which one will conduct most heat?\ntag{D}{215}{keph203}{14}{}{Class XI Part 2 \textperiodcentered{} Ch 10}}
{\opt{$r = 1\unit{cm}$, $l = 1\unit{m}$}{$r = 1\unit{cm}$, $l = \frac12\unit{m}$}{$r = 2\unit{cm}$, $l = 2\unit{m}$}{$r = 2\unit{cm}$, $l = \frac12\unit{m}$}}

\QQ{A Carnot engine working between $300\unit{K}$ and $400\unit{K}$ has $800\unit{J}$ of useful work. The amount of heat energy supplied to the engine from the source is\ntag{D}{238}{keph204}{13}{}{Class XI Part 2 \textperiodcentered{} Ch 11}}
{\opt{$2400\unit{J}$}{$3200\unit{J}$}{$1200\unit{J}$}{$3600\unit{J}$}}

\QQ{A particle executing SHM has a maximum speed of $0.5\unit{m\,s^{-1}}$ and maximum acceleration of $1.0\unit{m\,s^{-2}}$. The angular frequency of oscillation is\ntag{D}{266}{keph206}{8}{}{Class XI Part 2 \textperiodcentered{} Ch 13}}
{\opt{$2\unit{rad\,s^{-1}}$}{$0.5\unit{rad\,s^{-1}}$}{$2\pi\unit{rad\,s^{-1}}$}{$0.5\pi\unit{rad\,s^{-1}}$}}

\QQ{A source of sound is moving with a velocity of $50\unit{m\,s^{-1}}$ towards a stationary observer. The observer measures the frequency of sound as $500\unit{Hz}$. The apparent frequency of sound as heard by the observer when source is moving away from him with the same speed is (Speed of sound at room temperature $350\unit{m\,s^{-1}}$)\ntag{X}{}{}{}{}{Doppler effect}}
{\opt{$400\unit{Hz}$}{$666\unit{Hz}$}{$375\unit{Hz}$}{$177.5\unit{Hz}$}}

\QQ{If there are only one type of charge in the universe, then ($\vec E\to$ Electric field, $\vec{ds}\to$ Area vector)\ntag{D}{30}{leph101}{30}{ \textperiodcentered{} Exemplar 1.10}{Class XII Part 1 \textperiodcentered{} Ch 1}}
{\opt{$\oint\vec E\cdot\vec{ds}\ne0$ on any surface}{$\oint\vec E\cdot\vec{ds}$ could not be defined}{$\oint\vec E\cdot\vec{ds} = \infty$ if charge is inside}{$\oint\vec E\cdot\vec{ds} = 0$ if charge is outside, $= \dfrac{q}{\varepsilon_0}$ if charge is inside}}

\QQ{An electron of mass $m$, charge $e$ falls through a distance $h$ meter in a uniform electric field E. Then time of fall\ntag{D}{17}{leph101}{17}{ \textperiodcentered{} Example 1.7}{Class XII Part 1 \textperiodcentered{} Ch 1}}
{\opt{$t = \sqrt{\dfrac{2hm}{eE}}$}{$t = \dfrac{2hm}{eE}$}{$t = \sqrt{\dfrac{2eE}{hm}}$}{$t = \dfrac{2eE}{hm}$}}

\QQ{If $\vec E_{ax}$ and $\vec E_{eq}$ represents electric field at a point on the axial and equatorial line of a dipole. If points are at a distance $r$ from the centre of the dipole, for $r\gg a$\ntag{D}{24}{leph101}{24}{}{Class XII Part 1 \textperiodcentered{} Ch 1}}
{\opt{$\vec E_{ax} = \vec E_{eq}$}{$\vec E_{ax} = -\vec E_{eq}$}{$\vec E_{ax} = -2\vec E_{eq}$}{$\vec E_{eq} = 2\vec E_{ax}$}}

\QQ{Nature of equipotential surface for a point charge is\ntag{D}{54}{leph102}{10}{}{Class XII Part 1 \textperiodcentered{} Ch 2}}
{\opt{Ellipsoid with charge at foci.}{Sphere with charge at the centre of the sphere.}{Sphere with charge on the surface of the sphere.}{Plane with charge on the surface.}}

\QQ{A particle of mass $1\unit{g}$ and charge $1\unit{\mu C}$ is held at rest on a frictionless horizontal surface at distance $1\unit{m}$ from the fixed charge $2\unit{mC}$. If the particle is released, it will be repelled. The speed of the particle when it is at a distance of $10\unit{m}$ from the fixed charge\ntag{A}{56}{leph102}{12}{}{Class XII Part 1 \textperiodcentered{} Ch 2}\nsrc{None found in NEET/AIPMT/AIEEE/JEE papers, 2005--2026 (NCERT concept application)}}
{\opt{$60\unit{m\,s^{-1}}$}{$100\unit{m\,s^{-1}}$}{$90\unit{m\,s^{-1}}$}{$180\unit{m\,s^{-1}}$}}

\QQ{A capacitor of $8\unit{F}$ is connected as shown. Charge on the plates of the capacitor
\fig{\begin{circuitikz}[scale=0.6,transform shape,font=\Large]
\draw (0,4) to[C,l=$8\,\mathrm F$] (2.5,4) to[R,l=$20\,\Omega$] (5,4) -- (5,0);
\draw (0,4) -- (0,0);
\draw (0,2.2) to[R,l=$4\,\Omega$] (5,2.2);
\draw (0,0) -- (1.6,0) to[battery1,l_=$5\,\mathrm V$] (2.4,0) to[R,l_=$1\,\Omega$] (5,0);
\end{circuitikz}}\ntag{A}{72}{leph102}{28}{}{Class XII Part 1 \textperiodcentered{} Ch 2}\nsrc{None found in NEET/AIPMT/AIEEE/JEE papers, 2005--2026 (NCERT concept application)}}
{\opt{$32\unit{C}$}{$40\unit{C}$}{$0\unit{C}$}{$80\unit{C}$}}

\QQ{Four metal plates are arranged as shown. Capacitance between X and Y ($A\to$ Area of each plate, $d\to$ distance between the plates)
\fig{\begin{tikzpicture}[scale=0.8]
\draw[line width=1pt] (0,2.2) -- (2.4,2.2) node[pos=0.5,above,font=\tiny]{1}; \draw (2.4,2.2) -- (3.0,2.2) node[circ]{} node[right,lbl]{X};
\draw[line width=1pt] (0,1.5) -- (2.4,1.5) node[pos=0.5,above,font=\tiny]{2};
\draw[line width=1pt] (0,0.8) -- (2.4,0.8) node[pos=0.5,above,font=\tiny]{3}; \draw (2.4,0.8) -- (3.0,0.8) node[circ]{} node[right,lbl]{Y};
\draw[line width=1pt] (0,0.1) -- (2.4,0.1) node[pos=0.5,below,font=\tiny]{4};
\draw (0,1.5) arc (90:270:0.7);
\end{tikzpicture}}\ntag{A}{69}{leph102}{25}{}{Class XII Part 1 \textperiodcentered{} Ch 2}\nsrc{None found in NEET/AIPMT/AIEEE/JEE papers, 2005--2026 (NCERT concept application)}}
{\opt{$\dfrac32\dfrac{\varepsilon_0A}{d}$}{$\dfrac{2\varepsilon_0A}{d}$}{$\dfrac23\dfrac{\varepsilon_0A}{d}$}{$\dfrac{3\varepsilon_0A}{d}$}}

\QQ{Mobility of free electrons in a conductor is\ntag{D}{89}{leph103}{9}{}{Class XII Part 1 \textperiodcentered{} Ch 3}}
{\opt{directly proportional to electron density.}{directly proportional to relaxation time.}{inversely proportional to electron density.}{inversely proportional to relaxation time.}}

\QQ{Variation of resistance of the conductor with temperature is as shown
\fig{\begin{tikzpicture}[scale=0.8]
\draw[->] (-0.3,0) -- (3.4,0) node[right,lbl]{$T$}; \draw[->] (0,-0.2) -- (0,2.4) node[above,lbl]{$R$};
\node[lbl,below] at (0,-0.1) {$T_0$};
\draw (0,0.8) -- (2.8,2.3); \node[lbl] at (1.9,1.1) {Slope $= m$};
\draw[decorate,decoration={brace,amplitude=3pt}] (0.12,0.8) -- (0.12,0) node[midway,right=3pt,lbl]{$R_0$};
\end{tikzpicture}}
The temperature co-efficient ($\alpha$) of the conductor is\ntag{D}{90}{leph103}{10}{}{Class XII Part 1 \textperiodcentered{} Ch 3}}
{\opt{$\dfrac{R_0}{m}$}{$mR_0$}{$m^2R_0$}{$\dfrac{m}{R_0}$}}

\QQ{Potential difference between A and B in the following circuit
\fig{\begin{circuitikz}[scale=0.6,transform shape,font=\Large]
\draw (0,1.6) to[battery1,l_=$6\,\mathrm V$] (1,1.6) to[R,l_=$2\,\Omega$] (3,1.6) node[circ]{} node[above]{A} to[battery1,l_=$4\,\mathrm V$] (4.2,1.6) to[R,l_=$8\,\Omega$] (6.4,1.6) node[circ]{} node[above]{B};
\draw (0,1.6) -- (0,0) -- (6.4,0) -- (6.4,1.6);
\end{circuitikz}}\ntag{A}{97}{leph103}{17}{}{Class XII Part 1 \textperiodcentered{} Ch 3}\nsrc{None found in NEET/AIPMT/AIEEE/JEE papers, 2005--2026 (NCERT concept application)}}
{\opt{$4\unit{V}$}{$5.6\unit{V}$}{$2.8\unit{V}$}{$6\unit{V}$}}

\QQ{In the following network potential at `O'
\fig{\begin{circuitikz}[scale=0.6,transform shape,font=\Large]
\draw (0,0) node[circ]{} node[above]{A} node[below]{8 V} to[R,l_=$2\,\Omega$] (3,0) node[circ]{} node[above]{O};
\draw (3,0) to[R,l=$4\,\Omega$] (5,2) node[circ]{} node[right]{B, 4 V};
\draw (3,0) to[R,l_=$2\,\Omega$] (5,-2) node[circ]{} node[right]{C, 2 V};
\end{circuitikz}}\ntag{A}{97}{leph103}{17}{}{Class XII Part 1 \textperiodcentered{} Ch 3}\nsrc{None found in NEET/AIPMT/AIEEE/JEE papers, 2005--2026 (NCERT concept application)}}
{\opt{$4\unit{V}$}{$3\unit{V}$}{$6\unit{V}$}{$4.8\unit{V}$}}

\QQ{Effective resistance between A and B in the following circuit
\fig{\begin{circuitikz}[scale=0.55,transform shape,font=\Large]
\draw (0,3) to[R,l=$10\,\Omega$] (3,3) -- (6,3);
\draw (0,3) to[R,l_=$10\,\Omega$] (0,1.5) to[R,l_=$10\,\Omega$] (3,1.5) -- (4.6,1.5) node[circ]{} node[above]{A};
\draw (3,3) to[R,l=$10\,\Omega$] (3,1.5);
\draw (6,3) to[R,l=$10\,\Omega$] (6,1.5) -- (5.2,1.5) node[circ]{} node[above]{B};
\draw (0,1.5) -- (0,0) to[R,l_=$10\,\Omega$] (6,0) -- (6,1.5);
\end{circuitikz}}\ntag{A}{97}{leph103}{17}{}{Class XII Part 1 \textperiodcentered{} Ch 3}\nsrc{None found in NEET/AIPMT/AIEEE/JEE papers, 2005--2026 (NCERT concept application)}}
{\opt{$10\unit{\Omega}$}{$20\unit{\Omega}$}{$5\unit{\Omega}$}{$\dfrac{20}3\unit{\Omega}$}}

\QQ{Two heating coils of resistances $10\unit{\Omega}$ and $20\unit{\Omega}$ are connected in parallel and connected to a battery of emf $12\unit{V}$ and internal resistance $1\unit{\Omega}$. The power consumed by them are in the ratio\ntag{A}{93}{leph103}{13}{}{Class XII Part 1 \textperiodcentered{} Ch 3}\nsrc{None found in NEET/AIPMT/AIEEE/JEE papers, 2005--2026 (NCERT concept application)}}
{\opt{1:4}{1:3}{2:1}{4:1}}

\QQ{A proton is projected with a uniform velocity `$v$' along the axis of a current carrying solenoid, then\ntag{D}{109}{leph104}{3}{ \textperiodcentered{} Exemplar 4.4}{Class XII Part 1 \textperiodcentered{} Ch 4}}
{\opt{the proton will be accelerated along the axis.}{the proton path will be circular about the axis.}{the proton move along helical path.}{the proton will continue to move with velocity `$v$' along the axis.}}

\QQ{In the cyclotron, as radius of the circular path of the charged particle increases ($\omega$ = angular velocity, $v$ = linear velocity)\ntag{X}{}{}{}{}{cyclotron}}
{\opt{both $\omega$ and $v$ increases}{$\omega$ only increases, $v$ remains constant}{$v$ increases, $\omega$ remains constant}{$v$ increases, $\omega$ decreases}}

\QQ{A conducting wire carrying current is arranged as shown. The magnetic field at `O'
\fig{\begin{tikzpicture}[scale=0.85]
\coordinate (O) at (0,0);
\draw[->] (O) -- (30:1.6) node[pos=0.55,above,lbl]{$R_1$}; \draw[->] (30:2.8) -- (30:1.6);
\draw[->] (O) -- (-30:2.8); \draw (-30:2.8) -- (-30:1.6) node[pos=0.5,below,lbl]{$R_2$};
\draw[->] (-30:1.6) arc (-30:30:1.6); \draw[->] (30:2.8) arc (30:-30:2.8);
\draw (0.45,0) ++(-30:0) arc (0:0:0); \draw (-30:0.5) arc (-30:30:0.5); \node[lbl] at (0.8,0) {$60^\circ$};
\node[lbl,left] at (O) {O}; \node[lbl] at (30:3.0) {$I$};
\end{tikzpicture}}\ntag{A}{114}{leph104}{8}{}{Class XII Part 1 \textperiodcentered{} Ch 4}\nsrc{None found in NEET/AIPMT/AIEEE/JEE papers, 2005--2026 (NCERT concept application)}}
{\opt{$\dfrac{\mu_0i}{12}\left[\dfrac1{R_1}-\dfrac1{R_2}\right]$}{$\dfrac{\mu_0i}{12}\left[\dfrac1{R_1}+\dfrac1{R_2}\right]$}{$\dfrac{\mu_0i}{6}\left[\dfrac1{R_1}-\dfrac1{R_2}\right]$}{$\dfrac{\mu_0i}{6}\left[\dfrac1{R_1}+\dfrac1{R_2}\right]$}}

\QQ{The quantity of a charge that will be transferred by a current flow of $20\unit{A}$ over 1 hour 30 minutes period is\ntag{D}{82}{leph103}{2}{}{Class XII Part 1 \textperiodcentered{} Ch 3}}
{\opt{$10.8\times10^3\unit{C}$}{$10.8\times10^4\unit{C}$}{$5.4\times10^3\unit{C}$}{$1.8\times10^4\unit{C}$}}

\QQ{A galvanometer coil has a resistance of $50\unit{\Omega}$ and the meter shows full scale deflection for a current of $5\unit{mA}$. This galvanometer is converted into voltmeter of range $0$--$20\unit{V}$ by connecting\ntag{A}{131}{leph104}{25}{}{Class XII Part 1 \textperiodcentered{} Ch 4}\nsrc{None found in NEET/AIPMT/AIEEE/JEE papers, 2005--2026 (NCERT concept application)}}
{\opt{$3950\unit{\Omega}$ in series with galvanometer}{$4050\unit{\Omega}$ in series with galvanometer}{$3950\unit{\Omega}$ in parallel with galvanometer}{$4050\unit{\Omega}$ in parallel with galvanometer}}

\QQ{$\chi_1$ and $\chi_2$ are susceptibility of a paramagnetic material at temperatures $T_1$\,K and $T_2$\,K respectively, then\ntag{X}{}{}{}{}{Curie's law}}
{\opt{$\chi_1 = \chi_2$}{$\chi_1T_1 = \chi_2T_2$}{$\chi_1T_2 = \chi_2T_1$}{$\chi_1\sqrt{T_1} = \chi_2\sqrt{T_2}$}}

\QQ{At certain place, the horizontal component of earth's magnetic field is $3.0\unit{G}$ and the angle dip at that place is $30^\circ$. The magnetic field of earth at that location\ntag{X}{}{}{}{}{Earth's magnetism (angle of dip)}}
{\opt{$4.5\unit{G}$}{$5.1\unit{G}$}{$3.5\unit{G}$}{$6.0\unit{G}$}}

\QQ{The process of super imposing message signal on high frequency carrier wave is called\ntag{X}{}{}{}{}{communication systems (modulation)}}
{\opt{Amplification}{Demodulation}{Transmission}{Modulation}}

\QQ{A long solenoid with 40 turns per cm carries a current of $1\unit{A}$. The magnetic energy stored per unit volume is \rule{1cm}{0.4pt} $\mathrm{J/m^3}$.\ntag{A}{169}{leph106}{16}{}{Class XII Part 1 \textperiodcentered{} Ch 6}\nsrc{None found in NEET/AIPMT/AIEEE/JEE papers, 2005--2026 (NCERT concept application)}}
{\opt{$3.2\pi$}{$32\pi$}{$1.6\pi$}{$6.4\pi$}}

\QQ{A wheel with 10 spokes each of length `$L$' m is rotated with a uniform angular velocity `$\omega$' in a plane normal to the magnetic field `$B$'. The emf induced between the axle and the rim of the wheel.\ntag{D}{165}{leph106}{12}{ \textperiodcentered{} Example 6.7}{Class XII Part 1 \textperiodcentered{} Ch 6}}
{\opt{$\dfrac12N\omega BL^2$}{$\dfrac12\omega BL^2$}{$\omega bL^2$}{$N\omega BL^2$}}

\QQ{The rms value of current in a $50\unit{Hz}$ AC circuit is $6\unit{A}$. The average value of AC current over a cycle is\ntag{D}{178}{leph107}{2}{}{Class XII Part 1 \textperiodcentered{} Ch 7}}
{\opt{$6\sqrt2$}{$\dfrac{3}{\pi\sqrt2}$}{Zero}{$\dfrac{6}{\pi\sqrt2}$}}

\QQ{A capacitor of capacitance $10\unit{\mu F}$ is connected to an AC source and an AC Ammeter. If the source voltage varies as $V = 50\sqrt2\sin100t$, the reading of the ammeter is\ntag{D}{185}{leph107}{9}{}{Class XII Part 1 \textperiodcentered{} Ch 7}}
{\opt{$50\unit{mA}$}{$70.7\unit{mA}$}{$5.0\unit{mA}$}{$7.07\unit{mA}$}}

\QQ{In a series L.C.R circuit, the potential drop across L, C and R respectively are $40\unit{V}$, $120\unit{V}$ and $60\unit{V}$. Then the source voltage is\ntag{D}{187}{leph107}{11}{}{Class XII Part 1 \textperiodcentered{} Ch 7}}
{\opt{$220\unit{V}$}{$160\unit{V}$}{$180\unit{V}$}{$100\unit{V}$}}

\QQ{In a series L.C.R.\ circuit an alternating emf ($v$) and current ($i$) are given by the equation $v = v_0\sin\omega t$, $i = i_0\sin\left(\omega t+\dfrac\pi3\right)$. The average power dissipated in the circuit over a cycle of AC is\ntag{D}{191}{leph107}{15}{}{Class XII Part 1 \textperiodcentered{} Ch 7}}
{\opt{$\dfrac{v_0i_0}{2}$}{$\dfrac{v_0i_0}{4}$}{$\dfrac{\sqrt3}{2}v_0i_0$}{Zero}}

\QQ{Electromagnetic radiation used to sterilise milk is\ntag{D}{210}{leph108}{10}{}{Class XII Part 1 \textperiodcentered{} Ch 8}}
{\opt{X-ray}{$\gamma$-ray}{UV rays}{Radiowaves}}

\QQ{A plane glass plate is placed over a various coloured letters (violet, green, yellow, red). The letter which appears to raised more\ntag{A}{229}{leph201}{9}{}{Class XII Part 2 \textperiodcentered{} Ch 9}\nsrc{None found in NEET/AIPMT/AIEEE/JEE papers, 2005--2026 (NCERT concept application)}}
{\opt{Red}{Yellow}{Green}{Violet}}

\QQ{A ray of light passes through four transparent media with refractive index $n_1$, $n_2$, $n_3$ and $n_4$ as shown. The surfaces of all media are parallel
\fig{\begin{tikzpicture}[scale=0.7]
\foreach \x in {0,1,2,3,4} \draw (\x,0) -- (\x,3);
\draw (0,3) -- (4,3) (0,0) -- (4,0);
\foreach \x/\n in {0.5/1,1.5/2,2.5/3,3.5/4} \node[lbl] at (\x,2.7) {$n_\n$};
\draw (-0.2,0.5) -- (1,1.0) -- (2,1.15) -- (3,1.45) -- (4.2,1.95);
\node[lbl,left] at (-0.2,0.5) {A}; \node[lbl,above] at (1,1.0) {B}; \node[lbl,above] at (2,1.15) {C}; \node[lbl,above] at (3,1.45) {D}; \node[lbl,right] at (4.2,1.95) {E};
\end{tikzpicture}}
If the emergent ray DE is parallel to incident ray AB, then\ntag{A}{228}{leph201}{8}{}{Class XII Part 2 \textperiodcentered{} Ch 9}\nsrc{None found in NEET/AIPMT/AIEEE/JEE papers, 2005--2026 (NCERT concept application)}}
{\opt{$n_1 = n_4$}{$n_2 = n_4$}{$n_3 = n_4$}{$n_1 = \dfrac{n_2+n_3+n_4}{3}$}}

\QQ{Focal length of a convex lens is $20\unit{cm}$ and its RI is 1.5. It produces an erect, enlarged image if the distance of the object from the lens is\ntag{A}{235}{leph201}{15}{}{Class XII Part 2 \textperiodcentered{} Ch 9}\nsrc{None found in NEET/AIPMT/AIEEE/JEE papers, 2005--2026 (NCERT concept application)}}
{\opt{$40\unit{cm}$}{$30\unit{cm}$}{$15\unit{cm}$}{$20\unit{cm}$}}

\QQ{A ray of light suffers a minimum deviation when incident on an equilateral prism of refractive index $\sqrt2$. The angle of incidence is\ntag{A}{240}{leph201}{20}{}{Class XII Part 2 \textperiodcentered{} Ch 9}\nsrc{None found in NEET/AIPMT/AIEEE/JEE papers, 2005--2026 (NCERT concept application)}}
{\opt{$30^\circ$}{$45^\circ$}{$60^\circ$}{$50^\circ$}}

\QQ{In Young's double slit experiment the source is white light. One slit is covered with red filter and the other with blue filter. There shall be\ntag{A}{265}{leph202}{11}{}{Class XII Part 2 \textperiodcentered{} Ch 10}\nsrc{None found in NEET/AIPMT/AIEEE/JEE papers, 2005--2026 (NCERT concept application)}}
{\opt{Alternate red \& blue fringes}{Alternate dark \& pink fringes}{Alternate dark \& yellow fringes}{No interference}}

\QQ{Light of wavelength $600\unit{nm}$ is incident normally on a slit of width $0.2\unit{mm}$. The angular width of central maxima in the diffraction pattern is (measured from minimum to minimum)\ntag{D}{267}{leph202}{13}{}{Class XII Part 2 \textperiodcentered{} Ch 10}}
{\opt{$6\times10^{-3}\unit{rad}$}{$4\times10^{-3}\unit{rad}$}{$2.4\times10^{-3}\unit{rad}$}{$4.5\times10^{-3}\unit{rad}$}}

\QQ{For what distance is ray optics is good approximation when the aperture is $4\unit{mm}$ and the wavelength of light is $400\unit{nm}$?\ntag{X}{}{}{}{}{Fresnel distance}}
{\opt{$24\unit{m}$}{$40\unit{m}$}{$18\unit{m}$}{$30\unit{m}$}}

\QQ{The variation of photo-current with collector potential for different frequencies of incident radiation $\nu_1$, $\nu_2$ and $\nu_3$ is as shown in the graph, then
\fig{\begin{tikzpicture}[scale=0.8]
\draw[->] (-2.2,0) -- (2.8,0) node[below left,lbl]{Collector potential}; \draw[->] (0,0) -- (0,2.3) node[above,lbl]{Photo-current};
\draw (-0.6,0) .. controls (-0.2,0.9) and (-0.1,1.6) .. (0.1,1.62) -- (1.9,1.75);
\draw (-1.2,0) .. controls (-0.7,0.9) and (-0.3,1.6) .. (0.1,1.62);
\draw (-1.7,0) .. controls (-1.1,0.9) and (-0.5,1.6) .. (0.1,1.62);
\foreach \x/\n in {-0.6/01,-1.2/02,-1.7/03} \node[font=\tiny,below] at (\x,0) {$V_{\n}$};
\node[font=\tiny] at (-0.15,0.6) {$\nu_1$}; \node[font=\tiny] at (-0.75,0.8) {$\nu_2$}; \node[font=\tiny] at (-1.25,1.0) {$\nu_3$};
\draw[<->] (1.5,0.05) -- (1.5,1.7) node[midway,right,font=\tiny,align=left]{Saturation\\current};
\end{tikzpicture}}\ntag{D}{279}{leph203}{6}{}{Class XII Part 2 \textperiodcentered{} Ch 11}}
{\opt{$\nu_1 = \nu_2 = \nu_3$}{$\nu_1>\nu_2>\nu_3$}{$\nu_1<\nu_2<\nu_3$}{$\nu_3 = \dfrac{\nu_1+\nu_2}{2}$}}

\QQ{The de Broglie wavelength of an electron accelerated to a potential of $400\unit{V}$ is approximately\ntag{A}{285}{leph203}{12}{}{Class XII Part 2 \textperiodcentered{} Ch 11}\nsrc{None found in NEET/AIPMT/AIEEE/JEE papers, 2005--2026 (NCERT concept application)}}
{\opt{$0.03\unit{nm}$}{$0.04\unit{nm}$}{$0.12\unit{nm}$}{$0.06\unit{nm}$}}

\QQ{Total energy of electron in an excited state of hydrogen atom is $-3.4\unit{eV}$. The kinetic and potential energy of electron in this state\ntag{D}{305}{leph204}{16}{ \textperiodcentered{} Exercise 12.4}{Class XII Part 2 \textperiodcentered{} Ch 12}}
{\opt{$K = -3.4\unit{eV}$, $U = -6.8\unit{eV}$}{$K = 3.4\unit{eV}$, $U = -6.8\unit{eV}$}{$K = -6.8\unit{eV}$, $U = +3.4\unit{eV}$}{$K = +10.2\unit{eV}$, $U = -13.6\unit{eV}$}}

\QQ{When electron jumps from $n = 4$ level to $n = 1$ level, the angular momentum of electron changes by\ntag{D}{299}{leph204}{10}{}{Class XII Part 2 \textperiodcentered{} Ch 12}}
{\opt{$\dfrac{h}{2\pi}$}{$\dfrac{2h}{2\pi}$}{$\dfrac{3h}{2\pi}$}{$\dfrac{4h}{2\pi}$}}

\QQ{A radio-active sample of half-life 10 days contains $1000x$ nuclei. Number of original nuclei present after 5 days is\ntag{X}{}{}{}{}{radioactive decay law (half-life)}}
{\opt{$707x$}{$750x$}{$500x$}{$250x$}}

\QQ{An element X decays into element Z by two-step process.
\[\mathrm X\longrightarrow\mathrm Y+{}^4_2\mathrm{He},\qquad \mathrm Y\longrightarrow\mathrm Z+2\,\bar e\]
then\ntag{X}{}{}{}{}{alpha- and beta-decay}}
{\opt{X \& Z are isobars.}{X \& Y are isotopes.}{X \& Z are isotones.}{X \& Z are isotopes.}}

\QQ{A nucleus of mass $20\unit{u}$ emits a $\gamma$ photon of energy $6\unit{MeV}$. If the emission assume to occur when nucleus is free and rest, then the nucleus will have kinetic energy nearest to (take $1\unit{u} = 1.6\times10^{-27}\unit{kg}$)\ntag{A}{310}{leph205}{5}{}{Class XII Part 2 \textperiodcentered{} Ch 13}\nsrc{None found in NEET/AIPMT/AIEEE/JEE papers, 2005--2026 (NCERT concept application)}}
{\opt{$10\unit{keV}$}{$1\unit{keV}$}{$0.1\unit{keV}$}{$100\unit{keV}$}}

\QQ{Constant DC voltage is required from a variable AC voltage. Which of the following is correct order of operation?\ntag{X}{}{}{}{}{voltage regulation}}
{\opt{Regulator, filter, rectifier}{Rectifier, regulator, filter}{Rectifier, filter, regulator}{Filter, regulator, rectifier}}

\QQ{In a transistor, the collector current varies by $0.49\unit{mA}$ and emitter current varies by $0.50\unit{mA}$. Current gain $\beta$ measured is\ntag{X}{}{}{}{}{transistor}}
{\opt{49}{150}{99}{100}}

\QQ{Identify the logic operation carried out by the following circuit.
\fig{\begin{circuitikz}[scale=0.7,transform shape]
\draw (0,1.2) node[nand port](G1){}; \draw (0,-0.6) node[nand port](G2){}; \draw (2.6,0.3) node[nand port](G3){};
\draw (G1.in 1) -- ++(-0.3,0) coordinate(a1); \draw (G1.in 2) -- ++(-0.3,0) coordinate(a2); \draw (a1) -- (a2); \draw ($(a1)!0.5!(a2)$) -- ++(-0.5,0) node[circ]{} node[left]{A};
\draw (G2.in 1) -- ++(-0.3,0) coordinate(b1); \draw (G2.in 2) -- ++(-0.3,0) coordinate(b2); \draw (b1) -- (b2); \draw ($(b1)!0.5!(b2)$) -- ++(-0.5,0) node[circ]{} node[left]{B};
\draw (G1.out) -| ($(G3.in 1)+(-0.3,0)$) -- (G3.in 1); \draw (G2.out) -| ($(G3.in 2)+(-0.3,0)$) -- (G3.in 2);
\draw (G3.out) -- ++(0.4,0) node[circ]{} node[right]{Y};
\end{circuitikz}}\ntag{X}{}{}{}{}{digital electronics (logic gates)}}
{\opt{AND}{NAND}{NOR}{OR}}

%%BEGIN-ANALYSIS
\end{multicols}
\clearpage
\begin{center}{\pgfornament[width=5cm,color=gilt]{87}}\par
{\fontsize{20}{24}\selectfont\color{navy}\scshape KCET 2016 \textperiodcentered{} Paper Analysis}\par
{\color{gilt}\pgfornament[width=5cm,color=gilt,symmetry=h]{87}}\end{center}
\vspace{4pt}
\begin{multicols}{2}
\textbf{\color{crimson}1. Direct vs Applied}\par\smallskip
{\footnotesize Checked question by question against the four 2026-27 NCERT books (XI Part 1 \& 2, XII Part 1 \& 2). Out of 60, 48 questions are in the present syllabus; percentages are of these.}\par\smallskip
\begin{tabular}{@{}lrr@{}}\toprule
Direct from NCERT & \textbf{27} & 56\%\\
\quad text, summary, table or figure & 22 &\\
\quad NCERT solved Example & 4 &\\
\quad NCERT Exercise & 1 &\\
Applied (concept in NCERT) & \textbf{21} & 44\%\\\midrule
Analysed & 48 &\\
Deleted portion (not analysed) & 12 &\\\midrule
\textbf{Total} & \textbf{60} &\\\bottomrule\end{tabular}\par\smallskip
{\footnotesize Direct: Q 2, 3, 4, 8, 10, 11, 12, 13, 14, 16, 17, 18, 19, 23, 24, 29, 32, 38, 39, 40, 41, 42, 43, 49, 51, 53, 54.\\ From NCERT Examples: Q 4 (Example 4.7), 11 (Example 10.8), 17 (Example 1.7), 38 (Example 6.7).\\ From NCERT Exercises: Q 53 (Exercise 12.4).}\par\medskip
\textbf{\color{crimson}2. Where the applied questions resemble earlier exams}\par\smallskip
Of the 21 applied questions, \textbf{0} closely resemble a question from an earlier national exam:
NEET/AIPMT/AIIMS 0, JEE Main/Advanced 0, AIEEE 0. The other 21 (Q 1, 5, 6, 9, 20, 21, 22, 25, 26, 27, 28, 31, 33, 37, 44, 45, 46, 47, 48, 52, 57) had no close match in the NEET, AIPMT, AIIMS, AIEEE and JEE papers of 2005--2026 that were searched; they are straightforward applications of NCERT concepts.\par
\medskip
\textbf{\color{crimson}3. Difficulty level}\par\smallskip
\begin{tabular}{@{}lr@{}}\toprule
Easy & 32\\ Moderate & 16\\ Hard & 0\\\bottomrule\end{tabular}\par\smallskip
{\footnotesize Overall: \textbf{easy}; mostly single-formula numericals. The three circuit questions Q25--27 take the longest. Class XI: 13 questions, Class XII: 35 (analysed questions only). (Difficulty is the editor's assessment.)}\par\medskip
\columnbreak
\textbf{\color{crimson}4. Chapter-wise distribution}\par\smallskip
{\footnotesize\begin{tabular}{@{}p{3.6cm}c@{\,}r@{\,}r@{\,}r@{}}\toprule
Chapter & Cl. & Qs & Dir & App\\\midrule
Motion in a Straight Line & XI & 1 & 0 & 1\\
Motion in a Plane & XI & 2 & 2 & 0\\
Laws of Motion & XI & 2 & 1 & 1\\
Rotational Motion & XI & 1 & 0 & 1\\
Gravitation & XI & 1 & 1 & 0\\
Mech.\ Prop.\ of Solids & XI & 1 & 0 & 1\\
Mech.\ Prop.\ of Fluids & XI & 1 & 1 & 0\\
Thermal Prop.\ of Matter & XI & 2 & 2 & 0\\
Thermodynamics & XI & 1 & 1 & 0\\
Oscillations & XI & 1 & 1 & 0\\
Electric Charges \& Fields & XII & 3 & 3 & 0\\
Electrostatic Potential \& Capacitance & XII & 4 & 1 & 3\\
Current Electricity & XII & 7 & 3 & 4\\
Moving Charges \& Magnetism & XII & 3 & 1 & 2\\
Electromagnetic Induction & XII & 2 & 1 & 1\\
Alternating Current & XII & 4 & 4 & 0\\
Electromagnetic Waves & XII & 1 & 1 & 0\\
Ray Optics & XII & 4 & 0 & 4\\
Wave Optics & XII & 2 & 1 & 1\\
Dual Nature & XII & 2 & 1 & 1\\
Atoms & XII & 2 & 2 & 0\\
Nuclei & XII & 1 & 0 & 1\\
\bottomrule\end{tabular}}\par\medskip
\textbf{\color{crimson}5. Teaching takeaways}\par\smallskip
{\footnotesize\begin{itemize}[leftmargin=1em,nosep]
\item Solved Examples asked almost verbatim: Example 10.8 (pan cooling from $94^\circ$C to $86^\circ$C, p.\,221), Example 1.7 (time of fall of an electron, p.\,17), Example 6.7 (wheel with 10 spokes, p.\,165), Example 4.7 (box in an accelerating train, p.\,61); Exercise 12.4 (kinetic and potential energy from the total energy).
\item Most AC questions need only one textbook formula: $I=V\omega C$, $V=\sqrt{V_R^2+(V_C-V_L)^2}$, $P=\tfrac12v_0i_0\cos\phi$, average current zero over a cycle.
\item Twelve questions are on deleted portions (rolling, Doppler, cyclotron, Curie's law, dip, modulation, Fresnel distance, decay, regulation, transistor, logic gates); KCET 2016 was set on the pre-2023 syllabus.
\item Q33 (galvanometer to voltmeter) is a lab-manual question.
\end{itemize}}\par\medskip
\textbf{\color{crimson}6. Deleted portions}\par\smallskip
\textbf{12} of the 60 questions are on topics removed from the 2026-27 NCERT (rationalised syllabus). They are marked \textsc{Deleted Portion} and are not analysed: rolling motion (Q 7); Doppler effect (Q 15); cyclotron (Q 30); Curie's law (Q 34); Earth's magnetism (angle of dip) (Q 35); communication systems (modulation) (Q 36); Fresnel distance (Q 50); radioactive decay law (half-life) (Q 55); alpha- and beta-decay (Q 56); voltage regulation (Q 58); transistor (Q 59); digital electronics (logic gates) (Q 60).\par
\end{multicols}
%%END-ANALYSIS


